\chapter{The Shared Filter}
\label{ch:filter}

The device has one filter, shared by all the voices. It is a two-pole resonant filter of the state-variable kind, and
it can be a low-pass, a band-pass or a high-pass filter. Its input is the sum of the voices whose \texttt{FILTER\_EN}
bit is set; its output has its own volume and pan, and is mixed with the dry voices (Section~\ref{sec:filter-output}).
A voice that is filtered is not also heard dry.

\section{Operation}

The filter runs once for every output sample, all the time. It runs when no voice is routed to it, and while the output
is muted, so its state is always current. There is no register bit that turns it off. A program that does not want the
filter leaves its voices' \texttt{FILTER\_EN} bits clear.

The filter holds two state values, \emph{low} and \emph{band}. For each output sample, with the input \emph{in}, it
forms three results in this order:

\begin{align*}
	\mathrm{low}  & \leftarrow \mathrm{low} + f \times \mathrm{band}               \\
	\mathrm{high} & \leftarrow \mathrm{in} - \mathrm{low} - q \times \mathrm{band} \\
	\mathrm{band} & \leftarrow \mathrm{band} + f \times \mathrm{high}
\end{align*}

\noindent Each line uses the values the lines above it have just produced. The coefficient $f$ comes from the cutoff register
and $q$ from the resonance register (Sections~\ref{sec:filter-cutoff} and~\ref{sec:filter-resonance}). The output is
the new \emph{low}, \emph{band} or \emph{high}, as \texttt{FILTER\_CONFIG} selects: low-pass, band-pass or high-pass
(Table~\ref{tab:reg-filter}).

\begin{figure}[H]
	\centering
	\begin{tikzpicture}[
			scale=0.92, transform shape,
			box/.style={draw, rounded corners=2pt, fill=sirclightgray, align=center, font=\footnotesize, minimum height=0.9cm, minimum width=2.1cm},
			sum/.style={draw, circle, fill=white, font=\footnotesize, inner sep=1pt, minimum size=0.7cm},
			flow/.style={-{Stealth[length=2mm]}, thick},
			note/.style={font=\scriptsize, text=sircgray}
		]
		\node (in) at (-1.2,0) {in};
		\node[sum] (s) at (0.8,0) {$-$};
		\node[box] (f1) at (3.6,0) {$\times f$};
		\node[box] (b) at (6.6,0) {band\\integrator};
		\node[box] (f2) at (10.0,0) {$\times f$};
		\node[box] (l) at (13.0,0) {low\\integrator};
		\node[box] (q) at (5.0,-1.7) {$\times q$};
		\coordinate (bandtap) at (8.3,0);
		\coordinate (lowtap) at (14.7,0);
		\draw[flow] (in) -- (s);
		\draw[flow] (s) -- (f1);
		\draw[flow] (f1) -- (b);
		\draw (b) -- (bandtap);
		\draw[flow] (bandtap) -- (f2);
		\draw[flow] (f2) -- (l);
		\draw (l) -- (lowtap);
		\fill (bandtap) circle (0.06);
		\fill (lowtap) circle (0.06);
		\draw (bandtap) -- (8.3,-1.7);
		\draw[flow] (8.3,-1.7) -- (q.east);
		\draw (q.west) -- (0.8,-1.7);
		\draw[flow] (0.8,-1.7) -- (s.south);
		\draw (lowtap) -- (14.7,-3.0);
		\draw (14.7,-3.0) -- (1.2,-3.0);
		\draw[flow] (1.2,-3.0) |- (1.2,-0.5) -- (s.south east);
		\node[note, above] at (2.2,0.05) {high};
		\node[note, above] at (8.3,0.05) {band};
		\node[note, right] at (14.8,0) {low};
	\end{tikzpicture}
	\caption{The Filter}
	\label{fig:filter}
\end{figure}

\section{Arithmetic}

The filter works on signed 24-bit numbers with four fraction bits: the input sum of the voices' scaled samples is
multiplied by 16 to enter them, and an output is divided by 16 (with the shift rounding towards minus infinity) and
limited to 20 bits to leave them. The coefficients $f$ and $q$ are 8-bit numbers divided by 256, so each multiplication
is a state value times an 8-bit coefficient shifted right by 8. The shift truncates: it rounds towards minus infinity.
The filter needs three multiplications for each output sample: $f \times \mathrm{band}$, $q \times \mathrm{band}$ and
$f \times \mathrm{high}$.

Every sum and every state is limited to the 24-bit range and never wraps around. A loud input or a high resonance can
therefore drive the filter into the limit, and the states are held at the limit of their range: they do not wrap, so
the filter stays bounded. The filter does not reduce its gain to avoid this, and it does not remove the small errors
its truncation leaves: after a loud sound ends, the output settles to within a few output units of zero (at worst about
5), and does not fall further, provided that the resonance is below 255.

\section{Cutoff}
\label{sec:filter-cutoff}

The cutoff register \texttt{FILTER\_CUTOFF} holds an unsigned 8-bit value $c$, and $f = c / 256$. The nominal cutoff
frequency of the filter, which is the frequency of its poles when the damping is zero and is close to it for light
damping, is

\begin{equation*}
	f_c = \frac{f_s}{\pi} \arcsin\!\left(\frac{f}{2}\right)
\end{equation*}

\noindent where $f_s$ is the output sample rate of 31,960.226~Hz. The largest value, 255, gives about 5.3~kHz. With heavy damping the peak of the band-pass response lies a little above the nominal frequency, about 6\% above at the lowest resonance. The frequency is nearly
proportional to the register: each step is about 20~Hz at the low end, and a little more higher up.

\begin{table}[H]
	\centering
	\begin{tabular}{rrr}
		\toprule
		\textbf{Register} & \textbf{Frequency (Hz)} & \textbf{Next step up (semitones)} \\
		\midrule
		1                 & 19.9                    & 12.0                              \\
		2                 & 39.7                    & 7.0                               \\
		4                 & 79.5                    & 3.9                               \\
		8                 & 159.0                   & 2.0                               \\
		16                & 318.0                   & 1.05                              \\
		32                & 636.2                   & 0.53                              \\
		64                & 1275.0                  & 0.27                              \\
		128               & 2570.6                  & 0.14                              \\
		192               & 3910.6                  & 0.09                              \\
		254               & 5280.9                  & 0.07                              \\
		255               & 5303.8                  & ---                               \\
		\bottomrule
	\end{tabular}
	\caption{Nominal Filter Cutoff Frequencies}
	\label{tab:filter-cutoff}
\end{table}

Because the register is linear in frequency, the musical size of a step falls as the cutoff rises: the first step is an
octave and a step at 2.5~kHz is about a seventh of a semitone. A slow sweep through the low end is therefore heard as
steps.

A cutoff register value of 0 makes $f$ zero, and then neither of the two lines that change the state of the filter does
anything: \emph{low} and \emph{band} hold their last values. The low-pass output is a constant, and the high-pass
output is the input less that constant and the damping term. The filter resumes from the held state when the cutoff is
raised, with no step. A program that wants the filter silent sets the volume byte of \texttt{FILTER\_VOL\_PAN} to 0, or
clears \texttt{FILTER\_EN} on its voices.

\section{Resonance}
\label{sec:filter-resonance}

The resonance register \texttt{FILTER\_RESONANCE} holds an unsigned 8-bit value $R$. It sets the damping of the filter,
$q = d / 256$, where for $R$ from 0 to 254

\begin{equation*}
	d = 16 + \left\lfloor \frac{3 (255 - R)}{4} \right\rfloor
\end{equation*}

\noindent and for $R = 255$ the damping is zero. A smaller damping gives a sharper and taller peak at the pole frequency.

\begin{table}[H]
	\centering
	\begin{tabular}{rrrr}
		\toprule
		\textbf{Resonance} & \textbf{$d$ ($q \times 256$)} & \textbf{$Q = 1/q$} & \textbf{Band-pass peak (dB)} \\
		\midrule
		0                  & 207                           & 1.2                & 1.9                          \\
		64                 & 159                           & 1.6                & 4.2                          \\
		128                & 111                           & 2.3                & 7.3                          \\
		192                & 63                            & 4.1                & 12.2                         \\
		224                & 39                            & 6.6                & 16.4                         \\
		240                & 27                            & 9.5                & 19.6                         \\
		250                & 19                            & 13.5               & 22.7                         \\
		254                & 16                            & 16.0               & 24.2                         \\
		255                & 0                             & unbounded          & ---                          \\
		\bottomrule
	\end{tabular}
	\caption{Filter Resonance (Cutoff Register 64)}
	\label{tab:filter-resonance}
\end{table}

A resonance of 255 is a special value: the damping is zero, so the filter does not lose energy. Once an input or a
nonzero state sets it going, it oscillates at the pole frequency with a constant amplitude, and that amplitude is held
until the state is disturbed again. A filter whose states are zero and whose input is zero stays silent. Values near
255 are very resonant but do damp.

\section{Filter Output}
\label{sec:filter-output}

The selected output, divided by 16 and limited to 20 bits, is scaled by the filter volume, the high byte of
\texttt{FILTER\_VOL\_PAN}, and divided between the channels by the filter pan, its low byte, as a voice is
(Sections~\ref{sec:voice-gain} and~\ref{sec:mixing}). The two shares are added to the accumulators with the dry voices.
A modulation source can steer the cutoff and the filter pan (Chapter~\ref{ch:modulation}). The resonance and the volume
cannot be steered.

At reset the filter's state is zero (Section~\ref{sec:reset}).
